\slajdid{07-s019}
\begin{frame}[label=07-s019]
\gusttekst
Kada je reč o jezicima za opis hardvera, VHDL nije jedini. Verilog HDL postoji već dugi niz godina i ima veliki broj fanova koji ga koriste i za simulaciju i za sintezu. Kako se Verilog razlikuje od VHDL-a? Prvo, Verilog ima određenu prednost u dostupnosti simulacionih modela. Ubrzo nakon svog uvođenja (od strane Gatevai Design Automation sredinom 1980-ih), Verilog se uspostavio kao de fakto standardni jezik za ASIC simulacione biblioteke. Široka popularnost Verilog HDL-a za ovu svrhu ostavila je VHDL donekle iza u podršci za širok spektar ASIC uređaja. Inicijativa VITAL poboljšava VHDL standardnim metodom zapisa o kašnjenju koji je sličan onom koji se koristi u Verilog-u. VITAL će olakšati dobavljačima ASIC-a i drugima da brzo kreiraju VHDL modele iz postojećih Verilog biblioteka. Još jedna karakteristika koja je definisana u Verilogu (u Open Verilog International objavljenom standardu) je interfejs programskog jezika, ili PLI. PLI omogućava piscima simulacionih modela da izađu van Verilog-a kada je to potrebno za kreiranje bržih simulacionih modela, ili da kreiraju funkcije (koristeći jezik C) koje bi bilo teško ili neefikasno implementirati direktno u Verilogu. U korist VHDL-a su njegovo usvajanje kao IEEE standard (standard 1076) i njegove karakteristike upravljanja dizajnom višeg nivoa. Ove karakteristike upravljanja dizajnom, koje uključuju deklaraciju konfiguracije i karakteristike biblioteke VHDL-a, čine VHDL idealnim jezikom za upotrebu u velikim projektima. U većini aspekata, međutim, VHDL i Verilog su identični u funkciji. Sintaksa je drugačija (sa Verilogom koji liči na C, a VHDL više kao Pascal ili Ada), ali osnovni koncepti su isti.
\end{frame}
